\documentclass[12pt]{article}
\usepackage[utf8]{inputenc}
\usepackage[T1]{fontenc}
\usepackage{hyperref}
\title{Tawoyan Language Resources for K-12 Curriculum Development}
\author{Compiled by Claude Sonnet 5 (claude-sonnet-5) for Andr\'as Kornai}
\date{}
\begin{document}
\maketitle

\section{Introduction}
Tawoyan (ISO 639-3: twy; also Taboyan, Tabuyan, Tawoyan Dayak) is a Northeast Barito language of southeastern Kalimantan, Indonesia (Central and East Kalimantan). Sources found give lower speaker estimates than the 35,000 in the brief (one gives about 20,000). Search results turned up no dictionary, primer, textbook or school programme.

\section{Primary Educational Resources}
None found in search results.

\section{Linguistic and Reference Materials}
\begin{itemize}
\item Tawoyan (encyclopedic overview): \url{https://en.wikipedia.org/wiki/Tawoyan}
\item Northeast Barito languages: \url{https://en.wikipedia.org/wiki/Northeast_Barito_languages}
\item ASJP wordlist for Tawoyan: \url{https://asjp.clld.org/languages/TAWOYAN}
\item Numeral list (Channumerals database): \url{https://lingweb.eva.mpg.de/channumerals/Tawoyan.htm}
\item Glottolog map for Tawoyan: \url{https://glottolog.org/resource/languoid/id/tawo1243.bigmap.html}
\end{itemize}

\section{Religious and Community Resources}
Scripture Earth lists the language under ISO twy; I did not verify what, if any, scripture is available there: \url{https://www.scriptureearth.org/00eng.php?iso=twy}

\section{Digital and Audio Resources}
No audio, video or corpus resources were found.

\section{Curriculum Development Context}
Documentation is limited to word lists and classification. Curriculum work would require new documentation with speaker communities in Kalimantan.
\end{document}